\documentclass[11pt,a4paper]{article}

\usepackage[utf8]{inputenc}
\usepackage[cyr]{aeguill}
\usepackage[francais]{babel}

% Les packages
%\usepackage[french]{babel}
\usepackage{fancybox}
\usepackage{graphics}
\usepackage{color}
\usepackage{latexsym}
\usepackage{amsmath,amsthm,amssymb}
\usepackage[plain]{fullpage} 
\usepackage{verbatim}
\usepackage{times,epsfig}
\usepackage{listings}
\usepackage{multirow}
\usepackage{amsmath,amsthm,amssymb,amscd,txfonts,bbm} % RAJOUT ICI !!!
\usepackage{graphics,epsfig} % RAJOUT ICI !!!
%\usepackage{algorithmic}
\usepackage[final]{pdfpages} 
\usepackage{fancyvrb}
\usepackage{enumerate}
\usepackage{listingsutf8}
\usepackage{comment}
\usepackage{minted}

% small et verbatim: to be coherent with \file
\newenvironment{smallverbatim}%
{\endgraf\small\verbatim}{\endverbatim}

\newenvironment{qv}
{\quote\Verbatim}
{\endquote\endVerbatim}

\usepackage{lastpage}

\newcommand{\terme}[1]{\texttt{#1}}
\newcommand{\termSet}[1]{\{\texttt{#1}\}}
\newcommand{\guill}[1]{<<\,#1\,>>}
% Macro pour l'inclusion de figure  
\newcommand{\fig}[2]{%
  \begin{figure*}[hbt]
   \begin{center}
      \includegraphics[width=12cm]{#1}
  \caption{\label{#1}  #2}
  \end{center}
 \end{figure*}
}
\newcommand{\figTex}[2]{%
  \begin{figure*}[hbt]
 \centerline{\input{#1.pstex_t}}
  \caption{\label{#1}  #2}
 \end{figure*}
}

% Macro résumés
\definecolor{turquoise}{rgb}{0.20 ,0.60, 0.60}
\definecolor{vert}{rgb}{0.40,  0.60 ,0.00}
\definecolor{violet}{rgb}{0.80 , 0.00 , 0.80}
\definecolor{bleu}{rgb}{0.20,  0.20,  0.80}
\definecolor{rose}{rgb}{1.00,  0.20,  0.40}

\lstset{basicstyle=\footnotesize,showstringspaces=false,language=XML,breaklines=false,columns=fullflexible,frame=shadowbox, captionpos=b}

\newenvironment{solution}[0]{
~\\ \color{red} \textbf{Solution : }\color{black}}{%
}
\excludecomment{solution}


\newcommand{\cle}[1]{\textbf{#1}}
\def\fk#1{\textit{#1}}
\def\tble#1{\textit{#1}}

\parindent=0pt           
                             
\usepackage{url}

\newtheorem{Exercice}{Definition}


\begin{document}

\title{Sujet d'examen UE NFE204 FOAD - Session 1}
\author{Année universitaire 2019-2020, 30 juin}
\date{\empty}

\maketitle

{\bf \Large Consignes.}
Cet examen est exceptionnellement organisé à distance. 
Vous devez rendre votre copie numérisée à 21h00 au plus tard. 
Aucune copie ne sera acceptée après cette heure limite. 
Les seuls formats acceptés sont le PDF et le JPEG, ceci afin d'éviter tout problème de lecture.  
Prenez vos dispositions pour être en mesure de fournir le bon format le moment venu. Une activité 
de discussion en direct (\guill{\textit{chat}}) est ouverte dans l'espace Moodle pour d'éventuelles questions.

Important\,: il y a 22 points en tout à gagner. Lisez soigneusement l'énoncé et choisissez 
l'ordre selon vos points forts.

\section{Première partie : modélisation NoSQL et MapReduce (8 pts)}

En période d'épidémie, nous voulons construire un système de prévention. Ce système
doit être informé rapidement des nouvelles infections détectées, retrouver 
et informer rapidement les personnes ayant rencontré récemment une personne
infectée, et détecter enfin les foyers infectieux (\textit{clusters}). 
On s'appuie sur une application installée sur les téléphones mobiles. Elle fonctionne
par \textit{bluetooth}: quand deux personnes équipées de l'application sont proches 
l'une de l'autre, l'une d'entre elles (suivant un protocole d'accord)
 envoie au serveur un \emph{message 
de contact} dont voici  trois exemples\,:

\begin{small}
\begin{minted}{json}
{"_id": "7ytGy", "pseudo1": "xuyh57", "pseudo2": "jojoXYZ", "date": "30/06/2020"}
{"_id": "ui9xiuu", "pseudo1": "jojoXYZ", "pseudo2": "tat37HG", "date": "30/06/2020"}
{"_id": "Iuuu76", "pseudo1": "jojoXYZ", "pseudo2": "xuyh57", "date": "30/06/2020"}
\end{minted}
\end{small}

Ce message contient donc un identifiant unique, les pseudonymes
des deux personnes et la date de la rencontre. Un pseudonyme
est une clé chiffrée identifiant une unique personne. Pour des raisons
de sécurité, un nouveau pseudo est généré régulièrement et chaque application
conserve \emph{localement} la liste des pseudos engendrés. 
Les messages de contact (mais pas les listes de pseudos) sont stockés sur un serveur central
dans une base $DB_1$.
\vspace*{0.2cm}

\textbf{Questions}

\begin{enumerate}
  \item $DB_1$ est 
  asymétrique (un pseudo est stocké parfois dans le champ
    \texttt{pseudo1}, parfois dans le champ \texttt{pseudo2}) et  redondante
    puisque chaque contact apparaît plusieurs fois si deux personnes se sont
     rencontrés souvent dans la journée (cf. les exemples ci-dessus).    
     Proposez un traitement de type MapReduce qui produit, à partir
     de la base $DB_1$, une base $DB_2$ des \emph{rencontres quotidiennes}
     contenant un document par pseudo et par jour, avec la liste
    des contacts effectués ce jour-là. Pour les 
     exemples précédents, on devrait obtenir en ce qui concerne \texttt{jojoXYZ}\,:
     
\begin{small}
\begin{minted}{json}
{ "pseudo": "jojoXYZ",
  "date": "30/06/2020",
  "contacts": [{"pseudo": "tat37HG", "count": 1},
               {"pseudo": "xuyh57", "count": 2}]
}
\end{minted}
\end{small}
    
    Vous disposez d'une fonction $groupby()$ qui prend
    un ensemble de valeurs et produit
    une liste contenant chaque valeur et son nombre d'occurrences. 
    Par  exemple $groupby (x, y, x, z, y, x) = [(x,3), (y,2), (z,1)]$.
   Attention à l'asymétrie de représentation
    des pseudos dans les messages. 
    
    \item Voici quatre documents. les deux premiers sont stockés
    sur le serveur $S_1$, le troisième sur le serveur $S_2$, et le dernier sur le serveur $S_3$.
    \begin{itemize}
       \item \texttt{$d_1$: \{"\_id", "pseudo1": "X1", "pseudo2": "X2", "date": "30/06/2020" \}}
       \item \texttt{$d_2$: \{"\_id", "pseudo1": "X3", "pseudo2": "X2", "date": "30/06/2020" \}}
       \item \texttt{$d_3$: \{"\_id", "pseudo1": "X2", "pseudo2": "X4", "date": "30/06/2020" \}}
       \item \texttt{$d_4$: \{"\_id", "pseudo1": "X1", "pseudo2": "X4", "date": "30/06/2020" \}}
    \end{itemize}
         Inspirez-vous de la figure \url{http://b3d.bdpedia.fr/calculdistr.html#mr-execution-ex}
    pour montrer le déroulement du traitement MapReduce avec deux \textit{reducers}. 
     
    \item Initialement la base $DB_2$ est relationnelle\,: quel est son schéma pour pouvoir
    représenter correctement le contenu des documents des rencontres quotidiennes\,?
    
    \item Vérifiez que votre schéma est correct en donnant le contenu 
    des tables pour la rencontre du pseudo "jojoXYZ" le 30 juin (document ci-dessus).
 
    \item Maintenant $DB_2$ est une base NoSQL documentaire 
    et les rencontres sont stockées dans une collection \texttt{Rencontres}. 
     Quand une personne est infectée, elle transmet au
         serveur la liste de ses pseudos. Voici par exemple la liste
des pseudos conservés sur mon application mobile\,::

\begin{small}
\begin{minted}{json}
{"mesPseudos": ["jojoXYZ", "johj0N", "fjhukij87", "kodhvy"]}
\end{minted}
\end{small}

      Tous les jours on stocke les listes dans une collection \texttt{Infections}.
         Donnez la chaîne de traitement qui construit la collection de tous les pseudos qui ont
         été en contact avec une personne infectée. 
\end{enumerate}

\textbf{Important}\,: pour spécifier les chaînes de traitement, 
donner les fonctions de Map et de 
Reduce, en javascript, en Pig, 
en pseudo-code, au pire en langage naturel en étant le plus précis possible.

\begin{solution}
\begin{enumerate}
\item Appelons ce traitement ``Rencontres''. Voici la fonction de Map. La clé
est constituée de deux champs, et il faut bien penser à emettre deux messages
pour corriger l'asymétrie des messages de rencontre.
\begin{small}
\begin{minted}{javascript}
   function mapRencontre (doc)
   {
      emit ([doc.pseudo1, doc.date], "pseudo": doc.pseudo2)
      emit ([doc.pseudo2, doc.date], "pseudo": doc.pseudo1)
      }
   }
\end{minted}
\end{small}

  La fonction de Reduce se contente de renvoyer la
  liste obtenue après élémination des doublons.
  \begin{small}
\begin{minted}{javascript}
   function reduceRencontre ([pseudo, date], listeContacts)
   {
      return ([pseudo, date], groupby (listeContacts))
   }
\end{minted}
\end{small}
  
  \item Pour la base relationnelle, il faut une table des rencontres quotidiennes\,:
\begin{small}
\begin{minted}{sql}
   create table Rencontre (
      id: int not null,
      pseudo: int no null,
      dateRencontre: date not null,
      primary key (id)
      unique (pseudo, dateRencontre)
      )
\end{minted}
\end{small}
NB\,: on peut aussi prendre comme clé primaire la paire \texttt{(pseudo, dateRencontre)}
Et une table représentant les contacts 
\begin{small}
\begin{minted}{sql}
   create table Contacts (
      idRencontre: int not null,
      pseudo: int no null,
      count: varchar,
      primary key (idRencontre, pseudo),
      foreign key (idRencontre) references Rencontre
     )
\end{minted}
\end{small}

  
  \item Prenons Pig. Il faut commencer par un \texttt{flatten} sur les listes
  de pseudos pour obtenir une collection donnant tous les pseudos des personnes
  infectées.
\begin{small}
\begin{minted}{pig}
pseudosInfectés = foreach Infection generate "pseudoInfecté" , flatten(mesPseudos);
\end{minted}
\end{small}

   En effectuant la jointure entre  \texttt{pseudosInfectés} et \texttt{Rencontre}
   on obtient les contacts à prévenir.
   
\begin{small}
\begin{minted}{pig}
contactsAPrevenir = join Rencontre by "pseudo" , pseudosInfectés by "pseudoInfecté";
\end{minted}
\end{small}

\end{enumerate}

\end{solution}

\section{Deuxième partie : recherche d'information (5 pts)}


La base des rencontres peut être représentée
comme une matrice dans laquelle chaque vecteur horizontal représente les
rencontres effectuées par une personne dans le passé ave ses contacts.
Nous considérons la matrice suivante. Notez qu'elle est symétrique
et que nous plaçons sur la diagonale le nombre total de contacts 
effectués par une personne.

\begin{center}
\begin{tabular}{l|c|c|c|c}
            & tat37HG & xuyh57 & jojoXYZ & Ubbdyu \\ \hline  
   tat37HG  & $5$ &  $1$   & $1$     & $3$   \\
   xuyh57   & $1$    & $3$  & $2$   & $0$ \\
   jojoXYZ  & $1$   & $2$  &   $3$  &  $0$        \\
   Ubbdyu   &  $3$  &  $0$         & 0           & $3$     \\\hline
   \end{tabular}
\end{center}

On veut étudier les foyers infectieux en appliquant des méthodes vues en cours NFE204,
et une ébauche de classification $kMeans$. 
\begin{enumerate}
   \item Donnez les normes des vecteurs.
  \item On soupçonne que \texttt{jojoXYZ} et \texttt{Ubbdyu} sont à l'origine de 
  deux foyers $C_1$ et $C_2$. Donnez une mesure de distance basée
  sur la similarité cosinus entre ces deux pseudos, entre 
  \texttt{jojoXYZ} et \texttt{xuyh57} et finalement entre
   \texttt{Ubbdyu} et \texttt{tat37HG}.
  En déduire la composition des deux foyers.
  
  \item Outre la fonction cosinus, on dispose d'une fonction $centroid(G)$ qui calcule le centroïde
  d'un ensemble de vecteurs $G$.
  
  Quelle chaîne de traitement scalable permet de produire la classification
  de tous les pseudos dans l'un des deux foyers et d'obtenir
  les nouveaux centroïdes\,? 
  
  \item La chaîne de traitement précédente doit être répétée
     jusqu'à convergence pour obtenir un $kMeans$ complet. Si  on travaille
     avec Spark, quelles informations devraient être conservées dans
     un RDD persistant selon vous\,?
\end{enumerate}



\begin{solution}
\begin{enumerate}
\item Normes\,: $\sqrt{36}=6 ; \sqrt{14}=3,74 ; \sqrt{14}=3,74; \sqrt{18}=4,25$.

\item Calculons la similarité cosinus

    \begin{enumerate}
       \item $cos(\rm{tat37HG, xuyh57}) =\frac{5+3+2}{6 \times 3,74} = 0,44$
       \item $cos(\rm{tat37HG, jojoXYZ}) =\frac{5+2+3}{6 \times 3,74} = 0,44$
       \item $cos(\rm{tat37HG, Ubbdyu}) =\frac{15+9}{6 \times 4,25}=0,95$
       \item $cos(\rm{xuyh57, jojoXYZ}) =\frac{1+5+5}{14}=0,78$
       \item $cos(\rm{xuyh57, Ubbdyu}) =\frac{3}{3,74 \times 4,25}=0,18$
       \item $cos(\rm{jojoXYZ, Ubbdyu}) =\frac{3}{3,74 \times 4,25}=0,18$
    \end{enumerate}
    
    On distingue donc bien deux groupes: \texttt{jojoXYZ} est très proche 
    de \texttt{xuyh57}, et \texttt{tat37HG} est très proche 
    de \texttt{Ubbdyu}.
    
    \item La fonction de map calcule $cosinus()$ pour les deux centroïdes
    $C1.centre$ et  $C2.centre$ et affecte le vecteur au foyer le plus proche.
    
    \begin{small}
\begin{minted}{javascript}
   function mapKmeans (vecteur)
   {
      cos1 = cosinus (C1.centre, vecteur)
      cos2 = cosinus (C2.centre, vecteur)
      if ( cos1 > cos2) then
         emit ("C1", vecteur)
      else
         emit ("C2", vecteur)
      end;
      }
   }
\end{minted}
\end{small}
    
    La fonction de reduce calcule le centroide.
        \begin{small}
\begin{minted}{javascript}
   function reduceKmeans (key, L: liste(vecteurs))
   {
      return (key, centroid(L));
   }
\end{minted}
\end{small}
    
    \item La matrice doit être dans un RDD persistant. À chaque étape il  
    faut également conserver les centroïdes en mémoire RAM.
    
 \end{enumerate}
\end{solution}

\section{Troisième partie : analyse à grande échelle (5 pts)}

On veut maintenant contrôler la propagation du virus grâce aux informations
recueillies. On prend la  matrice des rencontres suivantes (la même
que précédemment mais
cette fois  la diagonale est à 0 puisqu'on s'intéresse à la transmission
d'une personne à l'autre).

\begin{center}
\begin{tabular}{l|c|c|c|c}
            & tat37HG & xuyh57 & jojoXYZ & Ubbdyu \\ \hline  
   tat37HG  & $0$ &  $1$   & $1$     & $3$   \\
   xuyh57   & $1$    & $0$  & $2$   & $0$ \\
   jojoXYZ  & $1$   & $2$  &   $0$  &  $0$        \\
   Ubbdyu   &  $3$  &  $0$         & 0           & $0$     \\\hline
   \end{tabular}
\end{center}

Faisons quelques hypothèses simplificatrices\,:
\begin{enumerate}
  \item le nombre de rencontres de chaque pseudo avec un contact
    représente la probabilité de rencontrer chacun des contacts. 
    Par exemple, le premier document ci-dessus nous dit que \texttt{jojoXYZ} 
    a deux fois plus de chances de rencontrer \texttt{xuyh57} que 
    de rencontrer \texttt{tat37HG}.
   \item La liste des  contacts  est fixe. 
    \item Chaque personne  fait exactement une rencontre par jour.
\end{enumerate}

En résumé, chaque personne rencontre exactement un de ses 3 contacts, chaque jour,
avec une probabilité proportionnelle à la fréquence antérieure des rencontres avec
ces mêmes contacts.

\begin{enumerate}
   \item Donnez la matrice des probabilités de contact obtenue à partir de la
   matrice des rencontres 

   \item Dessinez le graphe en étiquetant les arêtes par leur probabilité.
   \item On reçoit l'information que \texttt{jojoXYZ} est infecté\,: quelle
   est la probabilité que \texttt{Ubbdyu} le soit également deux jours plus tard\,?
   \item Si l'une des personnes de notre groupe est infectée,
     expliquez comment on calcule
     les probabilités d'infection de chaque personne $n$ jours plus tard. Donnez
     l'exemple de la première étape.
 \end{enumerate}

\begin{solution}

\begin{enumerate}

\item Voici la matrice de transition. La somme des termes sur chaque ligne doit être égale
   à 1.


\begin{center}
\begin{tabular}{l|c|c|c|c}
            & tat37HG & xuyh57 & jojoXYZ & Ubbdyu \\ \hline  
   tat37HG  & $0$ &  $1/5$   & $1/5$     & $3/5$   \\
   xuyh57   & $1/3$    & $0$  & $2/3$   & $0$ \\
   jojoXYZ  & $1/3$   & $2/3$  &   $0$  &  $0$        \\
   Ubbdyu   &  $1$  &  $0$         & 0           & $0$     \\\hline
   \end{tabular}
\end{center}


  \item Voir figure \ref{reseau-social}

\fig{reseau-social}{Réseau des contacts}

  \item On part avec le vecteur $(0, 0, 1, 0)$ représentant 
       le fait que \texttt{jojoXYZ} est infecté. On multiplie
       par chaque vecteur de transition et on obtient la
       probabilité $\frac{1}{3}$ d'infecter  \texttt{tat37HG},  $\frac{2}{3}$ 
       d'infecter \texttt{xuyh57}
      
      Seconde étape\,: le vecteur est $(1/3, 2/3, 0, 0)$. Si on part
       de \texttt{tat37HG}, ça nous laisse  une probabilité d'$\frac{3}{15}$
       d'infecter \texttt{Ubbdyu}, auquel il faut ajouter $\frac{2}{9}$
       si on part  de \texttt{xuyh57}. Donc : $\frac{19}{45}$ au final.
       
    \item C'est exactement PageRank: on multiplie le vecteur et la matrice
     jusqu'à convergence (garantie). 
\end{enumerate}
\end{solution}


\section*{Quatrième partie : réplication et partitionnement (4 pts)}

Donnez la solution de l'exercice \url{http://b3d.bdpedia.fr/sharding.html#ex-s3-1}
(Ajout d'un serveur avec hachage cohérent).
\end{document}
